\section{Related Work}
\label{sec:related}

\paragraph{Same score, different behaviour.} \citet{dutta2024accuracy} show that compressed models with matched accuracy flip up to 13.6\% of answers and argue for flip-based evaluation; \citet{singh2026certifying} find per-item churn about $5\times$ the net change and propose paired equivalence tests at a declared margin. Both work on static benchmarks and count changed answers. We work in a closed loop, where a changed action is not an answer but an intervention whose effect depends on what happens next, and we ask what each change causes rather than how many changes occur.

\paragraph{Quantized agents.} \emph{Flat Score, Amplified Failures} \citep{jang2026flat} is the closest study: BF16, FP8 and AWQ-INT4 agents on \tautwo{} retail and telecom with four trials per task, no aggregate change surviving correction, TOST at $\pm 7.5$ points, tool-name hallucinations up to $2.5\times$, and a tighter error budget re-exposing a gap. It compares independently run trajectories, so a failure cannot be traced to the quantized action that started it or separated from what the agent did afterwards; we fork identical histories to make that separation. Benchmarks of compressed agents \citep{dong2025acbench} and studies of task-dependent precision needs \citep{lu2026mixquant,zhang2026quantclaw} report scores, not consequences of individual actions. \emph{Dynamic Mixed-Precision Routing} \citep{li2026dynamicmp} reports that interaction steps differ in their sensitivity to precision and learns step-level routing between a high- and a low-precision model. Our one-replacement gain $G$ is the quantity such a router would harvest by sending one decision to the reference model; we measure it causally at the state where the two models first disagree, rather than learning a proxy for it.

\paragraph{Detecting that a model changed.} Model Equality Testing \citep{gao2025model} detects, from samples, whether an API serves the claimed model; AgentProv \citep{wang2026agentprov} audits agentic API providers with tool-use policy probes and detects silent quantization. Detectability of a difference is a different property from its consequence: a quantized agent can be perfectly distinguishable and still cause no task loss, or cause loss through a handful of actions. We take the difference as given and measure its consequences.

\paragraph{Forking agent trajectories.} \emph{The Replay Gap} \citep{gonuguntla2026replay} evaluates model switching in SWE agents by forking live trajectories, with same-model control forks as a reference; a model swap rewrites most post-fork actions, and control-fork divergence depends on the serving format (FP8 controls diverge in over 90\% of cases, AWQ controls stay near-identical). We share the forking idea but differ in the intervention and the estimand: we fork at the first quantization-induced action divergence, swap the current action and the continuation controller separately, and score the task outcome, so the change in success splits into current-action, continuation and interaction terms. Our serving is byte-deterministic (\S\ref{sec:conditions}), so a same-checkpoint fork is required to reproduce its original exactly; it is an acceptance test (\S\ref{sec:replay}) rather than a noise reference.

\paragraph{Noise in agent evaluation.} Agent scores vary across identical-setting runs: on \tautwo{}, single runs must differ by 7--8 (retail) or 10--11 (airline) points to exceed noise under sampling \citep{tau2issue540}; pass@1 moves 2.2--6.0 points at temperature 0 on SWE-bench Verified \citep{bjarnason2026randomness}; inconsistency across repeated trials has been quantified with intraclass correlation \citep{mustahsan2025stochasticity}, and cost-controlled, reproducible agent benchmarking has been advocated \citep{kapoor2024agentsmatter}; prompt perturbations inflate paired variance far beyond reruns on single-turn function calling \citep{chen2026noisefloor}; and benchmark construction itself can create false successes \citep{zhu2025abc}. Serving numerics are a known source of greedy-decoding divergence \citep{du2026greedy,yuan2025nondeterminism,he2025defeating}. These works motivate our choices to remove serving noise, to treat prompt nonces as a researcher-defined perturbation whose independent draws form replicate blocks, and to keep conclusions about a fixed task pool separate from conclusions about new tasks.
